% Topic Diffusion extension: do talking points recur across districts?
% Example phrases are real 8-grams from the 2018--2023 cleaned transcripts (derived_screened/transcripts/), 2026-11-09.
% District/state counts are over full transcripts (not yet restricted to public comment).

\section{Topic Diffusion Extension}\label{sec:diffusion}

We ask whether public commenters in different districts use phrasing too similar to have arisen
independently, and where such shared \emph{talking points} spread. Notation follows
\cref{sec:m-bow}: documents $d_n$, districts $i(n)$, periods $t(n)$.

\subsection{Approach 1: similarity-based detection}

\begin{enumerate}
  \item \textbf{Public comment and keyword contexts.} Limit documents to public-comment segments,
  with each commenter's turn as one document. For each keyword $j$ in our lexicon (e.g.\
  \emph{indoctrination}, \emph{mask}, \emph{parental rights}), keep the set $\mathcal{G}_j$ of
  8-grams containing $j$. Real examples for $j = $ ``indoctrination'':
  \begin{itemize}
    \item ``will have opened the floodgates of an indoctrination''
    \item ``public schools are being turned into indoctrination centers''
    \item ``it's game over the indoctrination must stop now''
    \item ``the process of teaching a person or group''
  \end{itemize}
  The last is a dictionary definition, recited in 13 districts in 10 states by both critics and
  defenders of the curriculum. ``Indoctrination centers'' appears in 16 districts in 14 states.
  The most widely shared 8-gram, ``to indoctrinate or persuade students to a particular [point of
  view]'' (16 districts, 7 states), is board-policy text read aloud. Restricting to public comment
  removes it.

  \item \textbf{Pairwise similarity.} For each keyword $j$, compute $S_j(g, g')$ for all pairs
  $g, g' \in \mathcal{G}_j$ from different districts, using one of:
  (a) fuzzy lexical matching (shingle overlap or edit distance, which tolerates caption errors);
  (b) cosine similarity of sentence embeddings; or
  (c) an LLM judging whether two contexts make the same point, applied to candidate pairs from (a)--(b).

  \item \textbf{Shared talking points.} Link $g$ and $g'$ when $S_j(g, g') > \tau$, with $\tau$
  tuned on a placebo set (e.g.\ 2018--19 comments). A cluster of linked 8-grams is a
  \emph{shared talking point} if it spans at least two districts in different counties.

  \item \textbf{Analysis.} For each talking point: its timing (first appearance, peak), its source
  (official guidance, news, advocacy materials, or none found), and its spread across district
  pairs by partisanship, distance and state.
\end{enumerate}

\subsection{Approach 2: tracking known phrases}

Some talking points are already known. We collect them from the verbatim illustrative quotes
returned by our GABRIEL scoring (\cref{sec:m-llm}), and from advocacy-group
materials and news coverage. Examples are ``parents' bill of rights'' and ``unmask our children.''
For each phrase, including caption variants (``parents bill of rights,'' ``parental bill of
rights''), we record each district's first mention and its share of meetings mentioning the
phrase by period. We then study adoption: how quickly phrases spread, and whether districts adopt
sooner after nearby or co-partisan districts do.

\subsection{Trial results}\label{sec:diffusion-results}
% Prototype run 2026-11-09 on the screened 2018--2023 sample (80,063 transcripts, 1,485 districts).
% Code: source/analysis/diffusion/approach{1,2}/; tables: output/tables/diffusion/approach{1,2}/.

\paragraph{Approach 1.} We ran steps 1--3 for nine keywords (\cref{tab:diffusion-kw}). Two spans
are linked if their 15-word windows have a fuzzy-match score of at least 80, or if they share at
least four uncommon word 5-grams (about eight identical consecutive words). We set this threshold
on spans around the neutral word \emph{students}: only 12 of its 26.3 million cross-district pairs
pass, so we expect about 14 false links per keyword at the 8,000-hit cap. A language model
(\texttt{gpt-6-luna}) labelled 3,267 hits for speaker role, quotation source and stance. On 38
hand-coded hits it agreed with us 92\% of the time on whether the speaker is a public commenter.
A rule-based public-comment flag reached 0.76--0.93 precision but only 0.57--0.79 recall against
these labels. Total API cost was \$0.50.

\begin{table}[H]
\centering\footnotesize
\caption{Shared talking points by keyword (Approach 1).}\label{tab:diffusion-kw}
\setlength{\tabcolsep}{4pt}
\resizebox{\linewidth}{!}{%
\begin{tabular}{lrrrrrrrl}
\toprule
Keyword & Hits & Districts & PC share & Links (exp.\ false) & TPs & PC TPs & Official & Largest PC TP \\
 & & & (LLM) & & & & share & (districts/states) \\
\midrule
indoctrinat* & 4,431 & 575 & 0.85 & 503 (4) & 33 & 19 & 0.21 & 2/2 \\
mask(s) & 286,327$^{*}$ & 1,355 & 0.45 & 95 (15) & 47 & 12 & 0.51 & 13/11 \\
critical race theory / CRT & 17,058$^{*}$ & 763 & 0.81 & 667 (14) & 134 & 50 & 0.33 & 33/20 \\
parental / parents' rights & 5,413 & 616 & 0.63 & 368 (7) & 72 & 31 & 0.42 & 11/9 \\
groom* & 3,872 & 724 & 0.48 & 173 (3) & 31 & 3 & 0.61 & 3/2 \\
pornograph* & 6,302 & 584 & 0.72 & 389 (9) & 44 & 8 & 0.43 & 4/3 \\
equity & 187,604$^{*}$ & 1,295 & 0.22 & 20 (15) & 15 & 1 & 0.60 & 2/2 \\
vaccin* & 174,960$^{*}$ & 1,266 & 0.39 & 51 (14) & 33 & 7 & 0.58 & 4/4 \\
transgender* & 11,868$^{*}$ & 708 & 0.64 & 1,197 (14) & 113 & 31 & 0.33 & 9/5 \\
\bottomrule
\end{tabular}}
\par\smallskip
\begin{minipage}{0.97\linewidth}\scriptsize
$^{*}$ Random sample of 8,000 hits used for pairwise matching. PC share: share of 150 random hits
whose speaker the LLM labels a public commenter. Links: cross-district pairs above the threshold;
expected false links from the \emph{students} placebo. TP: talking point, a cluster spanning at
least two districts in different counties. PC TPs: clusters whose hits are mostly public commenters.
Official share: share of TPs that are mostly official text, board members or staff.
\end{minipage}
\end{table}

For rare, contested keywords, links far exceed the placebo, and public-comment talking points
recur across states:
\begin{itemize}
  \item ``[we have] been told on multiple occasions that critical race theory is not being taught
  in our schools'' (33 districts, 20 states, 2021--23);
  \item a recited definition, ``critical race theory is an academic concept\ldots'' (18 districts,
  15 states);
  \item ``educating and empowering parents to defend their parental rights at all levels of
  government'' (11 districts, 9 states, 2021--23), an advocacy group's mission statement read aloud;
  \item ``I don't want to wear a mask, I don't want my kids to wear [one]'' (13 districts, 11
  states), which may be an everyday phrase rather than a shared script.
\end{itemize}
Many of the largest clusters are not public comment. They include statute or policy text (a Florida
obscene-materials statute, state parental-rights laws, ``grooming'' in abuse-prevention policy) and
quoted statistics or definitions. For frequent administrative keywords (\emph{equity},
\emph{vaccine}, \emph{mask}), the number of links is close to the placebo and most clusters are
official. The LLM judged pairs just below the threshold to be the same talking point 12--78\% of
the time, compared with 3.5\% for placebo pairs. The lexical rule therefore misses many
paraphrases.

\paragraph{Approach 2.} We tracked 34 phrases with caption-spelling variants: hand seeds, n-grams
nearly absent in 2018--19 but used in at least 20 districts and 5 states in 2020--23, and phrases
from the GABRIEL quotes. Among the 529 districts with meetings in both periods, median first use
was 2020-07 for \emph{anti-racist}, 2021-06 for \emph{critical race theory} and \emph{mask mandate},
2021-08 for \emph{parental rights}, 2022-06 for \emph{parents bill of rights} and 2022-09 for
\emph{Moms for Liberty}. Controlling for meetings posted and locale, a one-SD higher county GOP
share raises the probability of adopting right-coded phrases by 5--7 points and lowers it by 6--11
points for left-coded phrases.

We fit a linear probability model of first adoption, with phrase-by-month fixed effects, on
523,255 district--phrase--months with 4,946 adoptions and a base hazard of 0.95\% per month.
A one-SD higher share of districts within 100\,km that have already adopted raises the hazard by
about 17\% (coefficient 0.012, s.e.\ 0.004). For same-state districts the increase is about 40\%.
The 100\,km coefficient exceeds all 200 draws in which district locations are permuted
($p < 0.005$), and the within-state permutation gives $p = 0.01$. Shared regional shocks such as
state laws or mandates would produce the same pattern. Mentioning a phrase is also not endorsing it.

\paragraph{Assessment.} Measuring talking-point diffusion through text reuse is feasible for rare,
contested terms that the public brings to meetings. It does not work for frequent administrative
terms, where shared text is mostly official. Separating a shared script from a common everyday
phrase requires matching to source material (advocacy toolkits, statutes) or timing evidence from
Approach 2. Because the captions have no speaker labels, public-comment status rests on LLM labels.
